\label{sec:interleaved}

The original Sfumato design proposed a \emph{learned mode-router}
that, at every sub-block boundary, chooses among
\{extend-AR, diffuse-current-block, re-diffuse-last-K-blocks, branch,
commit\}. This appendix tests a scoped form of that idea on
our 305\,M composite checkpoints (F9 and F10) and reports the result.

\subsection{Scripted interleaving (Phase I.0)}

We chain \code{gen\_ar} and \code{diff\_revise} from
\code{e5/scripts/probe5\_mode\_switch.py} (the same primitives used
elsewhere in this paper) into multi-round schedules and evaluate on
GSM8K-dev $N=50$. Five configurations:

\begin{itemize}
  \setlength\itemsep{2pt}
  \item \code{single\_ar\_128} — pure AR (baseline).
  \item \code{single\_switch\_64\_32} — one round of AR(64) +
        diff-revise(32); the existing \code{probe5} mode switch.
  \item \code{interleaved\_16\_8\_x3} — AR(16) $\to$ diff(8), three
        rounds.
  \item \code{interleaved\_8\_4\_x6} — finer-grained alternation, six
        rounds.
  \item \code{interleaved\_32\_16\_x2} — coarser, two long rounds.
\end{itemize}

All configurations use the Tier-0.5 anti-rep settings (temperature
$0.8$, top-$p$ $0.9$, repetition penalty $1.15$ on the AR head,
count-based repetition penalty on the diff head). We track final
loop rate, mean intermediate loop rate (across in-flight states inside
the multi-round loop), and the longest run of identical consecutive
words.

\paragraph{F9 (FineWeb-Edu only, 100\,\% trained).}
\begin{center}
\begin{tabular}{lrrrr}
\toprule
config & acc & loop$_\text{final}$ & loop$_\text{intermediate}$ & max word run \\
\midrule
\code{single\_ar\_128} & 0\% & 0\% & 0\% & 2 \\
\code{single\_switch\_64\_32} & 0\% & 24\% & 12\% & 17 \\
\code{interleaved\_16\_8\_x3} & 0\% & 18\% & 12\% & 9 \\
\code{interleaved\_8\_4\_x6} & \textbf{2\%} & \textbf{2\%} & \textbf{1\%} & \textbf{6} \\
\code{interleaved\_32\_16\_x2} & 0\% & 38\% & 22\% & 10 \\
\bottomrule
\end{tabular}
\end{center}

On F9, fine-grained alternation
(\code{interleaved\_8\_4\_x6}) beats single-switch on every measured
axis: $+2$\,pp accuracy (the first non-zero hit on F9 in our entire
study), $-22$\,pp loop rate, $-11$ max word run. Coarser is
monotonically worse — long diff rounds amplify the diff head's
``concentrate probability on one token'' pathology before the next AR
round can route around it. We confirmed (Phase I.0 plan gate) that
intermediate loop rate does \emph{not} compound across cycles at fine
grain: $1$\,\% intermediate vs. $2$\,\% final means iteration
\emph{removes} loops at this configuration, not adds them.

\paragraph{F10 step 10\,000 (Q/A mix, 5.5\,\% trained).}
\begin{center}
\begin{tabular}{lrrrr}
\toprule
config & acc & loop$_\text{final}$ & loop$_\text{intermediate}$ & max word run \\
\midrule
\code{single\_ar\_128} & 2\% & 0\% & 0\% & 4 \\
\code{single\_switch\_64\_32} & \textbf{4\%} & \textbf{0\%} & \textbf{0\%} & \textbf{3} \\
\code{interleaved\_16\_8\_x3} & 0\% & 0\% & 0\% & 3 \\
\code{interleaved\_8\_4\_x6} & 0\% & 0\% & 0\% & 3 \\
\code{interleaved\_32\_16\_x2} & \textbf{4\%} & \textbf{0\%} & \textbf{0\%} & \textbf{3} \\
\bottomrule
\end{tabular}
\end{center}

The picture inverts. On F10, even at only $5.5\,\%$ of total training,
all configurations have $0\,\%$ loop rate, and the single-round
\code{single\_switch\_64\_32} or the coarse
\code{interleaved\_32\_16\_x2} reach $4\,\%$ accuracy while
fine-grained iteration ties with the lowest. The 5\,\% GSM8K Q/A mix
fixed the diff-head loop pathology directly. Iteration becomes
unnecessary — there is nothing for it to recover from.

\subsection{Heuristic router (Phase I.1)}

We then implemented three non-learned routers that pick the next
chunk's mode from the model's own signals:

\begin{itemize}
  \setlength\itemsep{2pt}
  \item \textbf{H1 entropy gate.} Switch to diff if the last-position
        AR-logit entropy exceeds $T = 4.5$\,nats.
  \item \textbf{H2 diversity gate.} Switch to diff if the unique token
        count in the last 8 generated tokens drops below $4$.
  \item \textbf{H3 diff-confidence gate.} After each AR chunk, enter a
        diff phase and keep iterating until mean top-1 probability
        across the just-unmasked positions exceeds $P = 0.5$.
\end{itemize}

\begin{center}
\begin{tabular}{lrrrr}
\toprule
heuristic & F9 acc & F9 loop & F10@10k acc & F10@10k loop \\
\midrule
H1 entropy $\ge 4.5$ & 0\% & 0\% & 2\% & 2\% \\
H2 diversity $< 4$ & 0\% & 0\% & 0\% & 0\% \\
H3 diff-conf $\ge 0.5$ & 0\% & 24\% & 0\% & 0\% \\
\midrule
best fixed (I.0) & \textbf{2\%} & \textbf{2\%} & \textbf{4\%} & \textbf{0\%} \\
\bottomrule
\end{tabular}
\end{center}

No heuristic beats the best fixed schedule on either substrate. H2
never triggers diff (Tier-0.5 sampling keeps token diversity high
enough that the gate stays AR-only); H1 fires too rarely (15--23\,\%
of chunks); H3 fires too often (75\,\% of chunks) and on F9
reproduces the single-switch baseline's loop pathology.

\subsection{Decision and scope claim}

Following the Phase I plan's decision rule (``if heuristics tie or
lose, learning a router is unlikely to extract more signal''), we
\emph{did not} launch the REINFORCE-trained router (Phase I.2) or the
joint backbone+router training (Phase I.3). The honest finding is:

\begin{quote}
\textbf{Iterative AR$\leftrightarrow$diff routing helps only when the
diff head is broken. When training data matches the eval distribution
(F10's Q/A mix), single-round AR-then-diff is enough; routing was a
band-aid over an undertrained-diff-head pathology that proper training
removes.}
\end{quote}

This is a scope-limit on the Sfumato learned-router thesis, not a
universal refutation. Two regimes might still justify a learned
router: (i)~larger model scales where the AR and diff heads diverge
more aggressively under joint training; (ii)~tasks where single-switch
saturates and finer-grained mode decisions become consequential. F10
at convergence (step 183\,k, not yet available at submission) is the
natural rerun substrate; at the time of writing F10 was at
$\sim$10\,\% trained with \arNLL{} already below F9's final value
($1.94$ vs.\ $3.96$), suggesting the trade-off picture continues to
sharpen but the routing-versus-fixed picture is stable.
